\section{Conclusion and Outlook}

In this paper, we first showed the existence of a class of algorithmic detectors based on LLMs that can successfully identify language cues of deception without the presence of other visual or audio cues.  We contribute a novel dataset \textsc{T4Text} for deception detection in the presence of objective truth and achieve a model performance comparable to human performance. We further find that our best model performs well in cases where humans perform poorly and discover novel language models that could augment human reasoning to detect deception, opening up the possibility of human-LLM collaborations to combat misinformation.  

This paper advocates for human-AI collaboration, emphasizing the need for additional evidence on human dependence on algorithms in detecting textual deception. Using the methodology of incentivized decision-making from behavioral economics, our ongoing work seeks to observe and categorically elucidate human reliance on our AI model in different conditions (e.g., black box, glass box). Obtaining predictive features from our best model and those that drive human performance can provide us with a comprehensive understanding of why humans perform poorly in text-based deception detection tasks and whether AI tools can aid humans in such a crucial task that has lasting social, economic, and ecological impacts.

% Our models correctly identifies lying in situations where humans fail, demonstrating its capacity to unearth novel avenues of reasoning that might assist humans in complex deception detection tasks. It has the potential of triggering novel research on not only textual but multimodal deception data, specifically for speech and the dialogue dimension, which we plan to be explore in the future.

\section{Limitations}
We acknowledge that online misinformation can be very different in nature than lies in our dataset; however, we find examples of false information in Quora, Reddit where non-experts with propaganda use strategies like half-truths to misguide people. \textsc{T4Text} is a relatively small dataset; however, we showed that statistically significant analysis can be done with it. The human prediction data is derived from the original game show; hence the setup may not match exactly when we are evaluating text models. We are running additional human experiments to gather true human performance on \textsc{\textsc{T4Text}}. 

\section{Ethical Concerns} 
The dataset is in English. The original sessions occurred in the 1950s; hence we do not observe an equitable diversity in gender when it comes to the gender of the contestants. For all sessions, there were two female judges and two male judges. We occasionally observe judges asking questions that are biased toward gender or race; hence, any model that will be trained on this dataset may risk containing similar bias. In our paper, we do not train any generative model on this data, minimizing that risk. We acknowledge the potential misuse of such truth-detection systems, and we are following up with controlled experiments to understand if humans would over-rely on such systems.

\paragraph{Full declaration on dataset use.} This work was done under \texttt{IRB\_00167477}, which was approved for using freely available ``To Tell The Truth'' videos for transcription (we use open-source models) and evaluating humans' and algorithms' ability to detect deception from text.
``To Tell The Truth'' videos used in this paper were produced by CBS from 1956-59 and retrieved from YouTube. We consider using YouTube videos for research purposes to fall under the ``fair use'' clause, as stated: \url{https://www.youtube.com/intl/en-GB/yt/about/copyright/fair-use/}. We do not use any other data except videos, especially anything that is produced by YouTube.
We are additionally inspired by \citet{soldner2019box}, who use similar YouTube videos (of a different gameshow) for detecting deception from multimodal cues.
We release the video transcripts, the only data used in our paper, under a Creative Commons license (CC BY 4.0 DEED).

\section*{Acknowledgements} We sincerely thank Haimanti Bhattacharya and Subhasish Dugar for improving the foundational idea of this work. We also thank Chris Callison-Burch, Tushar Khot, and Peter Clark from the Allen Institute for AI for their generous feedback on the conclusions presented in this work. 